% !TeX program = lualatex
% Kurzlebenslauf – kompilieren mit: lualatex shortcv.tex  (oder: latexmk -lualatex shortcv.tex)

\documentclass[11pt,a4paper]{article}

\usepackage[ngerman]{babel}
\usepackage{fontspec}
\setmonofont{DejaVu Sans Mono}[Scale=MatchLowercase]
\usepackage[a4paper,margin=2cm]{geometry}
\usepackage{fancyvrb}
\usepackage[hidelinks]{hyperref}

\hypersetup{
  pdftitle   = {Kurzlebenslauf -- Anton Rechenauer},
  pdfauthor  = {Anton Rechenauer},
  pdfsubject = {Lebenslauf},
  pdflang    = {de-DE}}

\pagestyle{empty}

\begin{document}
\begin{Verbatim}[fontsize=\small]
ein kurzer Lebenslauf – ja, das ist Angeberei

**** COMMODORE BASIC ****  15359 BYTES FREE  READY.
es beginnt 1982 mit einem PET 2001 in der Schule – drei Jahre Erfahrung!
1985 in der Arbeit ein Olivetti M24 und GW-BASIC, dazu eine CP/M-Z80-Maschine
zu Hause natürlich ein C64, ein Atari ST (mein erstes C-Programm), ein Amiga

nach 1990 ein paar Vorlesungen an der Uni
Fortran, Pascal, REXX, Unix, SQL, C/C++, Windows MFC, schließlich Java
ich erinnere mich an Gopher und Netscape, an einer Sun-SPARC-Workstation

nach 2000 beginnt meine IT-Karriere mit 8 Jahren Softwareentwicklung
Linux, Apache, MySQL, PHP – was sonst
HTML und CSS, JavaScript, VB6 und Access, VBA
weiter zu .NET, MS SQL Server, Crystal Reports
weiter zu Java Servlets, EJB, Tomcat, Postgres
irgendwann muss jeder reguläre Ausdrücke lernen – und wie man vi beendet

nach 2010, und einer kurzen Pause, reißt die Flut nicht mehr ab
immer öfter lerne ich nur so viel, dass ich die Fachdomäne verstehe
Android, SAP, MS Project, Primavera P6, SharePoint-Sites,
die AWS- und Google-Cloud, Embedded-Programmierung, DOORS, Jira,
Requirements Engineering, Angular, OpenGL, GIS, Perforce, SAP Fiori

nach 2020, eher zufällig, endlich Infrastruktur
Cisco- und Aruba-CLI, Switch- und Firewall-Administration,
Server-Hardware, GBICs, Patchfelder, iLOs und iDRACs,
RAID und SCSI, Storage-Appliances, VMware und vSphere,
Veeam, WSUS, Monitoring, VLANs und Spanning Trees,
Windows Active Directory, ein Blick in die IT/OT/PLS-Welt

nebenbei habe ich in meiner Freizeit 79 MOOCs (mit Zertifikat) absolviert,
über alles, was mich interessiert hat: Scala, Haskell, Erlang, R,
Machine Learning, Apache Spark und Hadoop, TensorFlow ...

jetzt ist 2026 – ich habe fast alles wieder vergessen.

manche kennen die meisten der hier stehenden Wörter nicht (die Glücklichen).

gegeben seien n Zahlen h_1, ..., h_n, ihr arithmetisches Mittel ist
\bar{h} = (h_1 + ... + h_n)/n. Eines der Mysterien der Statistik ist das
Problem, die Ausgangsmenge in Töpfe aufzuteilen: Sind die h_i etwa die
Stundensätze bestimmter IT-Fachleute, so kann man gewisse Teilmengen h'
betrachten und völlig unterschiedliche Mittelwerte \bar{h'} finden.
Kühne Vermutung: Wer das seit über 30 Jahren macht, liegt mit
seinem Angebot über dem Durchschnitt. So ist das nun mal.
\end{Verbatim}
\end{document}
